% ── §4  Design ────────────────────────────────────────────────────────────────
\section{XTAL Design}
\label{sec:design}

Figure~\ref{fig:arch} shows the \xtal architecture.
The system consists of four components: a Telemetry Collector, an
XGBoost Contextual Bandit Agent, a Safety Layer, and an Actuator.

\begin{figure}[t]
  \centering
  % Replace with actual architecture figure after experiments
  \fbox{\parbox{0.9\columnwidth}{\centering\vspace{2em}
    [Architecture diagram: Figure 1 from experiments/figures/]\vspace{2em}}}
  \caption{XTAL architecture. The 10-second control loop spans two system
    layers (Flink execution layer and Kafka ingest layer).}
  \label{fig:arch}
\end{figure}

\subsection{State Space}
\xtal observes a six-dimensional state vector $\mathbf{s}_t \in \mathbb{R}^6$
sampled every control interval:

\begin{center}
\small
\begin{tabular}{ll}
\toprule
Feature & Source \\
\midrule
$p99\_latency\_ms$ & Prometheus histogram quantile (Flink) \\
$p999\_latency\_ms$ & Prometheus histogram quantile (Flink) \\
$input\_buffer\_usage$ & \texttt{flink\_taskmanager\_job\_task\_buffers\_inPoolUsage} \\
$checkpoint\_duration\_ms$ & \texttt{flink\_jobmanager\_job\_lastCheckpointDuration} \\
$kafka\_consumer\_lag$ & \texttt{kafka\_consumergroup\_lag} (Kafka Exporter) \\
$records\_in\_per\_sec$ & Flink records-in rate (30\,s window) \\
\bottomrule
\end{tabular}
\end{center}

\subsection{Action Space}
\xtal selects from 14 discrete joint configurations
$(c, b, \ell) \in \mathcal{A}$, where $c$ is checkpoint interval (ms),
$b$ is buffer timeout (ms), and $\ell$ is Kafka linger (ms).
The space spans from the most aggressive
$(c=500, b=0, \ell=0)$ to the most conservative
$(c=10000, b=100, \ell=50)$, with intermediate points covering the
space identified as beneficial in the baseline study (\S\ref{sec:motivation}).
Discrete actions simplify the safety layer and make SHAP values directly
interpretable as ``why this configuration over another.''

\subsection{Reward Function}
The reward for the transition $(\mathbf{s}_t, a_t, \mathbf{s}_{t+1})$ is:

\begin{equation}
r_t = -\frac{\min(p99_{t+1},\; p99_{max})}{p99_{SLA}}
      \;+\; r_{SLA}
      \;+\; r_{tput}
      \;+\; r_{ckpt}
      \;+\; r_{bp}
\end{equation}

where $p99_{SLA}=500$\ms, $p99_{max}=5000$\ms,
$r_{SLA}=-20$ if $p99_{t+1}>p99_{SLA}$ else $0$,
$r_{tput}=0.5\,(\frac{tput_{t+1}}{tput_t}-1)$ (capped at $+0.25$),
$r_{ckpt}=-1$ if $ckpt_{t+1}>10{,}000$\ms,
$r_{bp}=-2$ if $buffer\_usage_{t+1}>0.9$.
The primary term penalises latency proportionally to SLA headroom.
The SLA penalty creates a hard cliff to avoid lingering near the boundary.

\subsection{XGBoost Contextual Bandit}
We maintain a single XGBoost regressor $f: \mathbb{R}^{9} \to \mathbb{R}$
with input $(\mathbf{s}_t \| a_t)$ (state\,||\,action features) and
output the predicted reward.
At each step:
\begin{enumerate}
  \item With probability $\epsilon$ (decaying $0.3\to0.05$ over 500 steps),
    select a random action (exploration).
  \item Otherwise, evaluate $f$ for all 14 actions and select
    $a^* = \arg\max_a f(\mathbf{s}_t \| a)$ (exploitation).
  \item Record the transition $(\mathbf{s}_t, a_t, r_t, \mathbf{s}_{t+1})$.
  \item Retrain $f$ on all collected experience every 20 steps.
\end{enumerate}

\subsection{Safety Layer}
Before any action reaches the Flink REST API or Kafka AdminClient,
the safety layer enforces hard constraints:
\begin{itemize}
  \item $500 \leq c \leq 30{,}000$\ms (never disable checkpoints;
    never checkpoint more than 2$\times$/s which would stall operators)
  \item $1 \leq p \leq 12$ (never stop processing; never exceed slot budget)
  \item $0 \leq b \leq 200$\ms (never introduce more than 200\ms buffering)
\end{itemize}
These bounds are enforced by \texttt{ControlAction.is\_safe()} before any
API call. An unsafe action is rejected and logged; the previous
configuration persists.

Additionally, when $p99 > 0.8 \times p99_{SLA}$ (approaching SLA boundary),
the predicted rewards of all actions with $c > 2000$\ms are penalised by
$-5.0$ before the argmax, biasing the agent toward low-overhead
configurations under pressure.

\subsection{SHAP Explainability}
After each exploitative decision, \xtal runs
\texttt{shap.TreeExplainer} on the chosen action's input row and records
the top-3 (feature, SHAP value) pairs in the decision log CSV.
A positive SHAP value for feature $x$ means $x$'s current value
\emph{increased} the predicted reward for the chosen action relative
to the expected output over the training set.
This lets an operator read, for example:

\begin{quote}\small
\textit{Action \#6 (ckpt=2000ms, buf=10ms, linger=5ms): \\
  Key drivers:\\
  \quad input\_buffer\_usage: $\downarrow$ decreases reward by 0.41\\
  \quad p99\_latency\_ms: $\uparrow$ increases reward by 0.28\\
  \quad checkpoint\_ms: $\uparrow$ increases reward by 0.10}
\end{quote}

That is: ``buffer backpressure is high (hurting us), but reducing
checkpoint overhead outweighs it.''
